% Nouveau EGA V (cote 25): archive cote 25, site page 61 (PDF page 62).
% typescript. Draft compared: none.
% AI-transcribed by Claude Opus 5.5 (claude-opus-5-5, adaptive thinking, effort high), Message Batches API, batch ega25b (msgbatch_01A7n5YuXubUPNK1zsepi8RG), 2026-11-13; not yet checked by eye.
\scanpage{25/61}
\mspage{56}

L'énoncé \unsure{10.2.} est une conséquence à peu près immédiate du suivant,
qui est tout à fait indépendant d'histoires de sections hyperplanes, et
serait sans doute mieux à sa place dans  un nº en rabiot aux
``immersions régulières''.

















Lorsque dans 10.2.  $S = \mathrm{Spec}\,k$ et X est de dimension $\unsure{\leqslant}1$,
alors l'hypothèse de régularité faite équivaut à $T = \emptyset$ , de sorte
que $Y_{\underline{L}}\to X$ est un \uline{isomorphisme} . On trouve ainsi\unsure{;} en conjuguant
avec 9.2. :

\uline{Corollaire} \unsure{10.4.} (à \unsure{10.2.}) Soit X une courbe lisse géom connexe dans
 l'espace projectif P sur un corps infini k, et soit $n\hand{\geqslant}2$.
Alors il existe un pinceau linéaire de n-formes sur P, définissant un
morphisme $X\to P^1$ ayant les propriétés suivantes: ce morphisme est

\mspage{61}
